\chapter{Grigory Margulis (2020)}

\section{Biographical Sketch}

Grigory Margulis was born on 24 February 1946 in Moscow, Russia. He studied at Moscow State University, where he completed his PhD in 1970 under the supervision of Yakov Sinai. Despite facing significant restrictions as a Jewish mathematician in the Soviet Union, Margulis produced work of extraordinary depth. He held positions at Moscow State University and the Institute for Information Transmission Problems before emigrating to the United States in 1986, joining Yale University. Margulis received the Fields Medal in 1978, the Wolf Prize in 2005, and the Abel Prize in 2020 (shared with Hillel Furstenberg).

\section{High-Level View for a General Audience}

\subsection{What Did Margulis Do?}

Imagine taking a flat sheet of paper and cutting out identical copies of the same tile, then laying them side by side to cover the entire plane. This repeating pattern is a discrete arrangement inside a continuous space. Margulis studied the deepest possible version of this question: when does a discrete repeating pattern inside a highly symmetric continuous space have only two possible behaviors---either it is rigidly structured (discrete) or it spreads out everywhere (dense)?

He discovered that in many important cases, there is no middle ground. This ``discrete or dense'' dichotomy became the engine behind some of the most profound results in modern mathematics, from the structure of lattices in Lie groups to the Oppenheim conjecture about quadratic forms.

\subsection{Why Does It Matter?}

Margulis proved that the discrete subgroups of semisimple Lie groups---the mathematical objects that encode symmetry---are astonishingly rigid. His superrigidity theorem shows that these subgroups cannot be deformed in the way one might na\"ively expect. This rigidity has far-reaching consequences: it solves the Oppenheim conjecture (a 60-year-old problem about quadratic forms), it classifies all lattices in higher-rank Lie groups, and it provides the foundation for the study of hyperbolic manifolds via the Margulis lemma. His work, alongside Hillel Furstenberg with whom he shared the 2020 Abel Prize, established homogeneous dynamics as a central pillar of modern mathematics.

\section{The Origin of the Problem}

The study of homogeneous spaces $G/\Gamma$ (where $G$ is a Lie group\index{Lie groups} and $\Gamma$ is a lattice) lies at the intersection of geometry, group theory, and dynamics.

\begin{knowledgebridge}{What Is a Homogeneous Space?}
Imagine the sphere $S^2$ and the group $SO(3)$ of all its rotations. The sphere is \emph{homogeneous} because any two points can be connected by a rotation. A homogeneous space $G/\Gamma$ generalizes this: $G$ is a continuous symmetry group and $\Gamma$ a discrete subgroup that ``cuts'' the space into repeating fundamental domains. These spaces are the natural stage for studying the interplay of geometry, group theory, and dynamics.
\end{knowledgebridge}

\subsection{Early History}

The problem of understanding the distribution of orbits in homogeneous spaces dates back to the work of Gauss on binary quadratic forms and the reduction theory of lattices. In the 20th century, the study of flows on homogeneous spaces became central to number theory through the work of Hedlund, Artin, and Siegel.

\subsection{Margulis' Starting Point}

Margulis began with the theory of discrete subgroups of Lie groups\index{Lie groups}, motivated by the work of Selberg, Borel, and Mostow on lattices in semisimple Lie groups\index{Lie groups}. The central questions were:
\begin{itemize}
\item Which groups admit lattices?
\item How are lattices classified?
\item What are the dynamics of group actions on homogeneous spaces?
\end{itemize}

\begin{primer}{What Is a Lattice in a Lie Group?}
A lattice $\Gamma$ in a Lie group $G$ is a discrete subgroup such that the quotient $G/\Gamma$ has finite volume. Think of $\mathbb{Z}^2$ inside $\mathbb{R}^2$: it is discrete, and the quotient $\mathbb{R}^2/\mathbb{Z}^2$ is a torus of finite area. In higher-rank Lie groups, lattices are surprisingly rigid---unlike in rank 1, they cannot be deformed continuously---and this rigidity is the engine behind Margulis' deepest results.
\end{primer}

\section{Deep Insight into the Problem}

\subsection{Margulis' Measure Rigidity}

Margulis realized that invariant measures for group actions on homogeneous spaces must satisfy strong rigidity properties. A measure $\mu$ on $G/\Gamma$ that is invariant under a subgroup $H \subseteq G$ cannot be ``spread out'' unless $H$ is large enough. This principle, refined by Ratner into a complete classification, is the key to understanding the behavior of unipotent flows.

\subsection{The Superrigidity Phenomenon}

Let $\Gamma$ be an irreducible lattice in $G = \prod_i G_i(k_i)$, a product of simple algebraic groups over local fields, with total rank at least 2. Margulis showed that any homomorphism $\Gamma \to H(k)$ (where $H$ is an algebraic group over a local field $k$) extends to a continuous homomorphism $G \to H(k)$, up to compact factors. This is called superrigidity.

The superrigidity theorem implies that every irreducible lattice in a higher-rank semisimple Lie group is arithmetic: it arises from the integer points of an algebraic group defined over a number field. This completely resolves the problem of classifying lattices in higher-rank groups, a problem that had seemed intractable.

\subsection{The Normal Subgroup Theorem}

For an irreducible lattice $\Gamma$ in a higher-rank semisimple Lie group $G$, any normal subgroup $N \trianglelefteq \Gamma$ is either:
\begin{itemize}
\item Contained in the center (hence finite), or
\item Of finite index in $\Gamma$.
\end{itemize}

This is remarkable: most groups have many normal subgroups (e.g., free groups have uncountably many), but higher-rank lattices are almost simple. The proof uses the dynamics of the action of $\Gamma$ on the homogeneous space $G/\Gamma$, showing that any non-central normal subgroup must act in a way that forces it to be large.

\subsection{The Oppenheim Conjecture}

Let $Q$ be an indefinite quadratic form in $n \geq 3$ variables that is not a scalar multiple of a rational form. The Oppenheim conjecture (1930) states that $Q(\mathbb{Z}^n)$ is dense in $\mathbb{R}$. Margulis proved this using unipotent flows on $SL(n,\mathbb{R})/SL(n,\mathbb{Z})$.

His insight: the values $Q(x)$ are related to the orbit of the identity under a one-parameter unipotent subgroup $u(t) \in SL(n,\mathbb{R})$:
\[
Q(x) = Q_0(g(t) \cdot \mathbb{Z}^n), \quad g(t) = \exp(tN),
\]
where $N$ is a nilpotent matrix chosen so that $u(t)$ preserves $Q_0$ but expands $\mathbb{Z}^n$. Then $Q(\mathbb{Z}^n)$ is the set of values of $Q_0$ at points in the orbit $u(t) \cdot \mathbb{Z}^n$. By Ratner's theorem, the closure of the orbit is itself a homogeneous space.

\begin{analogy}{Discrete or Dense---No Middle Ground}
Imagine throwing a dart at an infinite field. Either the dart lands on a marked grid point (discrete) or it lands arbitrarily close to any spot (dense). The Oppenheim conjecture says that for an indefinite quadratic form $Q$ in $n \geq 3$ variables, the set $Q(\mathbb{Z}^n)$ admits \textbf{no third option}: it must be dense in $\mathbb{R}$ unless it is obviously discrete for rational reasons. Margulis proved this by studying how unipotent flows move the integer lattice.
\end{analogy}

\subsection{The Margulis Lemma}

The Margulis lemma states that any discrete subgroup of $SO(n,1)$ generated by sufficiently small rotations is nilpotent. More precisely, there exists a constant $\varepsilon_n > 0$ (the Margulis constant) such that for any discrete subgroup $\Gamma \subseteq SO(n,1)$, the subgroup generated by elements with displacement less than $\varepsilon_n$ is nilpotent. This is a key ingredient in the study of hyperbolic manifolds: it gives the thick-thin decomposition, where the thin part consists of tubes around short geodesics.

\subsection{Margulis Numbers}

The Margulis constant for hyperbolic space: there exists $\varepsilon_n > 0$ such that in any hyperbolic $n$-manifold, every cyclic subgroup generated by a translation of length $<\varepsilon_n$ is central in the fundamental group. This constant plays a crucial role in the geometry of hyperbolic manifolds, bounding the size of the thin region.

\section{Biggest Contributions and New Tools}

\subsection{The Superrigidity Theorem}
Any irreducible lattice in a semisimple Lie group of rank at least 2 is arithmetic in the sense that it arises from the integer points of an algebraic group. This completely classified irreducible lattices in semisimple Lie groups.

\subsection{The Normal Subgroup Theorem}
Normal subgroups of irreducible higher-rank lattices are either finite or of finite index, proving that these lattices are almost simple.

\subsection{The Oppenheim Conjecture}
For an indefinite quadratic form in $n \geq 3$ variables that is not proportional to a rational form, the set $Q(\mathbb{Z}^n)$ is dense in $\mathbb{R}$. The proof introduced the use of unipotent flows to number theory.

\subsection{The Margulis Lemma}
Any discrete subgroup of $SO(n,1)$ generated by sufficiently small rotations is nilpotent. This is a key ingredient in the study of hyperbolic manifolds and their thick-thin decomposition.

\subsection{Margulis Numbers}
The Margulis constant for hyperbolic space: there exists $\varepsilon_n > 0$ such that in any hyperbolic $n$-manifold, every cyclic subgroup generated by a translation of length $<\varepsilon_n$ is central.

\subsection{Arithmeticity of Lattices}
Margulis completely classified the irreducible lattices in semisimple Lie groups, showing they are arithmetic when the group has rank at least 2.

\begin{whyitmatters}
Margulis' work created a highway between number theory, dynamics, and geometry. Before him, the distribution of integer solutions to equations seemed like a separate subject from the study of dynamical systems on homogeneous spaces. His ideas now permeate modern mathematics---from the proof of the Oppenheim conjecture to the classification of approximate groups and the study of expander graphs.
\end{whyitmatters}

\section{Impact of the Discovery in Science and Technology}

\subsection{Impact on Mathematics}

\begin{itemize}
\item \textbf{Number Theory}: The Oppenheim conjecture and its generalizations, the theory of quadratic forms, and the study of Diophantine approximation via homogeneous dynamics.
\item \textbf{Lie Theory}: The complete classification of lattices in semisimple Lie groups, the arithmeticity theorem, and the theory of discrete subgroups.
\item \textbf{Geometry}: The geometry of hyperbolic manifolds, the Margulis lemma, and the thick-thin decomposition.
\item \textbf{Dynamics}: Ratner's theorems on unipotent flows, which grew out of Margulis' work on the Oppenheim conjecture, completely classify invariant measures and orbit closures for unipotent flows on homogeneous spaces.
\end{itemize}

\subsection{Impact on Other Fields}

\begin{itemize}
\item \textbf{Mathematical Physics}: The study of quantum chaos and arithmetic quantum unique ergodicity, building on tools from homogeneous dynamics.
\item \textbf{Computer Science}: Expander graphs and the theory of pseudo-randomness, building on the work of Margulis on expander families.
\item \textbf{Engineering}: The theory of random walks on groups and their applications to networks and algorithms. Expander graphs constructed using Margulis' techniques are used in building high-throughput routing networks, error-correcting codes, and hash functions.
\end{itemize}

\section{Current Developments}

\subsection{The Work of Eskin--Mirzakhani}

Alex Eskin and Maryam Mirzakhani used Ratner's theorems and Margulis' ideas to classify invariant measures on moduli spaces of flat surfaces, proving the magic wand theorem.

\subsection{Arithmetic Quantum Unique Ergodicity}

Elon Lindenstrauss proved the arithmetic quantum unique ergodicity conjecture using tools from homogeneous dynamics, building on the framework Margulis helped create.

\subsection{Approximate Groups}

The structure theory of approximate groups, developed by Breuillard, Green, and Tao, uses ideas from the theory of Lie groups and homogeneous dynamics, tracing back to Margulis' work on lattices.

\subsection{Superstrong Approximation}

Bourgain, Gamburd, and Sarnak developed superstrong approximation, using expansion properties of Cayley graphs of arithmetic groups. This builds directly on Margulis' construction of expander families.

\subsection{Applications to Diophantine Approximation}

Recent work by Kleinbock, Margulis, and others on Diophantine approximation on manifolds uses homogeneous dynamics to prove results about the distribution of rational approximations to real numbers.

\section{Conclusion}

Grigory Margulis transformed the study of discrete subgroups of Lie groups, proving results of astonishing depth and generality. His superrigidity theorem and normal subgroup theorem completely changed our understanding of lattices in higher-rank groups. His resolution of the Oppenheim conjecture brought the powerful machinery of homogeneous dynamics to bear on number theory. Alongside Hillel Furstenberg, with whom he shared the 2020 Abel Prize, Margulis revealed that the dynamics of group actions on homogeneous spaces is a unifying theme connecting number theory, geometry, and combinatorics. The Abel Prize citation recognized Margulis (jointly with Furstenberg) ``for pioneering the use of methods from probability and dynamics in group theory, number theory and combinatorics.''

\section{Behind the Breakthrough: The Human Story}

Margulis won the Fields Medal in 1978 for his work on lattices in Lie groups, but the Soviet government denied him a visa to travel to Helsinki to receive it. The resulting international outcry among mathematicians highlighted the political restrictions faced by Jewish scientists in the Soviet Union. Despite these obstacles---and despite being denied permission to travel abroad for years---Margulis continued to produce work of breathtaking originality. He was finally permitted to emigrate to the United States in 1986, joining Yale University. His story is a testament to the resilience of mathematical creativity in the face of political oppression.

\section{Pedagogical Metaphors and Lineage}
\subsection{Signature Metaphor}
``Discrete subgroups of continuous symmetry groups that are so rigid they can only be either perfectly structured or everywhere dense---no middle ground.''

\subsection{Mathematical Lineage}
Margulis was advised by Yakov Sinai.

\section{Applied Science and Computational Algorithms}
\subsection{Key Takeaways for Applied Science}
Expander graphs constructed using Margulis' techniques are used in building high-throughput routing networks, error-correcting codes, and hash functions.

\subsection{Algorithms and Numerical Implementations}
Belief Propagation decoding algorithms for Low-Density Parity-Check (LDPC) codes operate on expander graphs built from arithmetic group lattices. Margulis' explicit constructions of expander families remain among the most practical.

\section{The Research Frontier}
Investigating random walks on general groups and lattices in Lie groups of rank 1; extending Ratner's theorems to more general settings; and establishing expander graph families from arithmetic subgroups of new types.

\section{Guided Exercises and Challenges}
\begin{enumerate}[label=\textbf{\arabic*.}]
  \item State Margulis' superrigidity theorem for lattices in higher-rank semisimple Lie groups and explain why it implies arithmeticity.
  \item Explain how the Oppenheim conjecture for indefinite quadratic forms in $n \geq 3$ variables follows from the study of unipotent flows on $SL(n,\mathbb{R})/SL(n,\mathbb{Z})$.
  \item Describe the Margulis lemma and its role in the thick-thin decomposition of hyperbolic manifolds.
\end{enumerate}

\section{Curriculum Map and Further Reading}
For students wishing to delve deeper into the mathematical machinery underlying these breakthroughs, the following learning path is recommended:
\begin{itemize}
  \item G. A. Margulis, \emph{Discrete Subgroups of Semisimple Lie Groups}, Springer.
  \item D. W. Morris, \emph{Introduction to Arithmetic Groups}, Deductive Press.
  \item D. Kleinbock and G. A. Margulis, \emph{Flows on Homogeneous Spaces and Diophantine Approximation}, Cambridge University Press.
\end{itemize}

\section*{Bibliography}
\begin{enumerate}
\item G. A. Margulis, ``Discrete subgroups of semisimple Lie groups,'' in \emph{Proceedings of the ICM 1978}, 149--154.
\item G. A. Margulis, ``Arithmeticity of irreducible lattices in semisimple groups of rank greater than 1,'' Appendix to \emph{Discrete Subgroups of Semisimple Lie Groups}, Springer, 1991.
\item G. A. Margulis, ``Oppenheim conjecture,'' in \emph{Fields Medalists' Lectures}, World Scientific, 1997.
\item G. A. Margulis, \emph{Discrete Subgroups of Semisimple Lie Groups}, Springer, 1991.
\item G. A. Margulis, ``Explicit constructions of expanders,'' Problems of Information Transmission 9 (1973), 325--332.
\item G. A. Margulis, ``Complete classification of discrete subgroups of semisimple Lie groups,'' in \emph{Proceedings of the ICM 1978}.
\item M. Ratner, ``On Raghunathan's measure conjecture,'' Annals of Mathematics 134 (1991), 545--607.
\item M. Ratner, ``Raghunathan's topological conjecture and distributions of unipotent flows,'' Duke Mathematical Journal 63 (1991), 235--280.
\item D. W. Morris, \emph{Introduction to Arithmetic Groups}, Deductive Press, 2015.
\item A. Eskin and M. Mirzakhani, ``Invariant measures on the space of flat surfaces,'' arXiv:1302.1003, 2013.
\item E. Lindenstrauss, ``Pointwise theorems for amenable groups,'' Inventiones Mathematicae 146 (2001), 259--295.
\item E. Lindenstrauss, ``Invariant measures and arithmetic quantum unique ergodicity,'' Annals of Mathematics 163 (2006), 165--219.
\item E. Breuillard, ``A brief introduction to approximate groups,'' in \emph{Thin Groups and Superstrong Approximation}, Cambridge University Press, 2014.
\item Z. Rudnick and P. Sarnak, ``The behaviour of eigenstates of arithmetic hyperbolic manifolds,'' Communications in Mathematical Physics 161 (1994), 195--213.
\item J. Bourgain and A. Gamburd, ``Uniform expansion bounds for Cayley graphs of $SL_2(\mathbb{F}_p)$,'' Annals of Mathematics 167 (2008), 625--642.
\item A. Eskin, ``Quasi-isometric rigidity of nonuniform lattices in higher-rank symmetric spaces,'' Journal of the American Mathematical Society 11 (1998), 321--363.
\item J. Bourgain, A. Gamburd, and P. Sarnak, ``Affine linear sieve, expanders, and sum-product,'' Inventiones Mathematicae 179 (2010), 559--644.
\item D. Kleinbock and G. A. Margulis, ``Flows on homogeneous spaces and Diophantine approximation on manifolds,'' Annals of Mathematics 148 (1998), 339--360.
\end{enumerate}
